%% maestro2tex -- generated figure; do not edit by hand.
\providecommand{\MaestroRoot}{include/analog/}
\begin{figure}[htbp]
  \centering
  {\pgfplotsset{compat=1.18}%
  \begin{tikzpicture}
    \begin{axis}[
      width=0.82\linewidth, height=6.2cm,
      xlabel={DAC code}, ylabel={Differential nonlinearity~[LSB]},
      grid=both,
      major grid style={line width=0.2pt, draw=black!14},
      minor grid style={line width=0.2pt, draw=black!7},
      minor tick num=1,
      axis line style={line width=0.4pt, draw=black!45},
      tick style={line width=0.4pt, draw=black!45},
      tick align=outside,
      label style={font=\small}, tick label style={font=\small},
      enlarge x limits=0.02, enlarge y limits=0.10,
      every axis plot/.append style={line join=round, no marks},
      legend style={font=\footnotesize, draw=none, fill=none,
                    at={(0.5,1.02)}, anchor=south, legend columns=-1,
                    /tikz/every even column/.append style={column sep=9pt}},
      legend cell align=left,
      scaled y ticks=false, scaled x ticks=false,
      y tick label style={/pgf/number format/fixed,
                          /pgf/number format/precision=2}
    ]
      \addplot[mtxA, line width=1.1pt, solid] table {\MaestroRoot data/dacr2r12_dnl_ideal_00.dat};
      \addlegendentry{tt}
      \addplot[mtxB, line width=1.1pt, dashed] table {\MaestroRoot data/dacr2r12_dnl_ideal_01.dat};
      \addlegendentry{ff}
      \addplot[mtxC, line width=1.1pt, dotted] table {\MaestroRoot data/dacr2r12_dnl_ideal_02.dat};
      \addlegendentry{fs}
      \addplot[mtxD, line width=1.1pt, dashdotted] table {\MaestroRoot data/dacr2r12_dnl_ideal_03.dat};
      \addlegendentry{sf}
      \addplot[mtxE, line width=1.1pt, densely dashed] table {\MaestroRoot data/dacr2r12_dnl_ideal_04.dat};
      \addlegendentry{ss}
    \end{axis}
  \end{tikzpicture}}
  \caption[{Differential nonlinearity against code over five process corners, conditions as\,\dots}]{Differential nonlinearity against code over five process corners, conditions as Figure~\ref{fig:dacr2r12-inl-ideal}. DNL is the first difference of the transfer divided by the endpoint LSB, minus one. The worst step is the mid-scale major carry at code 2048, \(-0.50\) to \(-0.54\,\mathrm{LSB}\); the sub-carries measure \(-0.27\) at code 1024, \(-0.23\) at 3072, \(-0.14\) at 512 and \(-0.13\,\mathrm{LSB}\) at 2560. The 4095-point curve is stride-decimated for the plot, so the comb drawn is a sample of the code grid and not every tooth of it; the extreme is preserved exactly. The converter is monotonic at every corner, with \(0.46\,\mathrm{LSB}\) of margin to the \(-1\,\mathrm{LSB}\) limit at worst.}
  \label{fig:dacr2r12-dnl-ideal}
\end{figure}
